\chapter{Risk Metrics}
\label{chap:risk_metrics}

\section{Overview}

Once a simulation has been run, the analyst typically wants to summarise
the loss distribution using a small number of risk metrics. \actsim
provides functions to compute the most widely used metrics directly from
simulated data.

This chapter defines each metric formally and describes how to obtain it
from the \texttt{StochasticSimulator} class.

\section{Mean and Standard Deviation}

The mean and standard deviation of the simulated aggregate losses are the
simplest summary statistics:
\[
    \bar{S} = \frac{1}{M} \sum_{m=1}^{M} S_m,
    \qquad
    \hat{\sigma}_S = \sqrt{\frac{1}{M-1} \sum_{m=1}^{M} (S_m - \bar{S})^{2}},
\]
where $S_m$ is the aggregate loss in simulation $m$ and $M$ is the
number of simulations. Both can be obtained directly from the
\texttt{results} attribute of the simulator:

\begin{lstlisting}[style=pythonstyle]
mean_loss = simulator.results.mean()
std_loss  = simulator.results.std()
\end{lstlisting}

\section{Percentiles}

A percentile of the simulated aggregate losses is computed using
\texttt{calc\_agg\_percentile}:

\begin{lstlisting}[style=pythonstyle]
p95 = simulator.calc_agg_percentile(95)
p99 = simulator.calc_agg_percentile(99)
\end{lstlisting}

\section{Value at Risk (VaR)}

The Value at Risk at confidence level $q \in (0, 1)$ is defined as
the $q$-th quantile of the aggregate loss distribution:
\[
    \mathrm{VaR}_q(S) = \inf \{ s : F_S(s) \geq q \},
\]
where $F_S$ is the cumulative distribution function of $S$. In a
simulation context, $\mathrm{VaR}_q$ is estimated by the empirical
$q$-quantile of the simulated values. The
\texttt{analyze\_results} method returns VaR at user-specified
quantiles:

\begin{lstlisting}[style=pythonstyle]
metrics = simulator.analyze_results(quantiles=[0.95, 0.99])
print(metrics["VaR"])
\end{lstlisting}

\section{Tail Value at Risk (TVaR)}

The Tail Value at Risk, also known as the conditional tail expectation
or expected shortfall, is the expected loss conditional on exceeding
the corresponding VaR threshold:
\[
    \mathrm{TVaR}_q(S) = \mathbb{E}\big[ S \mid S > \mathrm{VaR}_q(S) \big].
\]

In simulation, $\mathrm{TVaR}_q$ is estimated by averaging the
simulated values that exceed the empirical $\mathrm{VaR}_q$.

\begin{lstlisting}[style=pythonstyle]
metrics = simulator.analyze_results(quantiles=[0.95, 0.99])
print(metrics["TVaR"])
\end{lstlisting}

TVaR is a coherent risk measure and is preferred over VaR for
applications where the behaviour of the loss distribution beyond
the VaR threshold is material.

\section{Aggregate Exceedance Probability (AEP)}

The Aggregate Exceedance Probability at level $q$ is the $q$-th
percentile of the \emph{aggregate} (annual) loss distribution:
\[
    \mathrm{AEP}_q = \inf \{ s : F_S(s) \geq q \}.
\]
For a single-period simulation, AEP and VaR coincide. AEP is widely
reported in catastrophe risk and reinsurance contexts as a function of
return period, where return period $T$ corresponds to confidence level
$q = 1 - 1/T$.

\begin{lstlisting}[style=pythonstyle]
metrics = simulator.analyze_results(quantiles=[0.95, 0.99])
print(metrics["aep"])
\end{lstlisting}

\section{Occurrence Exceedance Probability (OEP)}

The Occurrence Exceedance Probability at level $q$ is the $q$-th
percentile of the \emph{maximum single-event} loss within a year. If
$M$ is the maximum single-event loss in a year,
\[
    \mathrm{OEP}_q = \inf \{ m : F_M(m) \geq q \}.
\]

OEP requires event-level data and therefore can only be computed when
the simulator is run with \texttt{keep\_all=True}. In that case
\texttt{analyze\_results} returns OEP alongside the other metrics:

\begin{lstlisting}[style=pythonstyle]
simulator = StochasticSimulator(
    "poisson", (10,),
    "lognormal", (10, 0.5),
    num_sim=10000,
    keep_all=True,
    seed=1234,
)
simulator.gen_agg_simulations()
metrics = simulator.analyze_results(quantiles=[0.95, 0.99])
print(metrics["oep"])
\end{lstlisting}

\section{Putting it Together}

A typical risk-metrics report contains VaR, TVaR, AEP, and OEP at a
small set of confidence levels. The example below produces such a
report and prints it as a single DataFrame.

\begin{lstlisting}[style=pythonstyle]
quantiles = [0.5, 0.75, 0.9, 0.95, 0.99, 0.995]
report = simulator.analyze_results(quantiles=quantiles)
print(report.round(2))
\end{lstlisting}

The output table contains one row per quantile and one column per
metric. This format is convenient for inclusion in technical memos and
audit packages.

\section{Choice of Metric}

Different metrics emphasise different aspects of the loss distribution:

\begin{itemize}[leftmargin=*]
    \item \textbf{VaR} is simple and widely understood, but does not
          describe behaviour beyond the threshold.
    \item \textbf{TVaR} captures tail behaviour and is preferred in
          frameworks that require coherent risk measures.
    \item \textbf{AEP and OEP} are particularly common in
          property-catastrophe and reinsurance contexts, where annual
          aggregate exposure and the largest single event need to be
          distinguished.
\end{itemize}

The choice of metric should be driven by the regulatory framework, the
contract structure, and the question being asked. \actsim makes all of
these metrics available from a single simulation run, allowing the
analyst to report whichever combination is most relevant.
